\documentclass[11pt]{amsart}

\usepackage[T1]{fontenc}
\usepackage{lmodern}
\usepackage{microtype}
\usepackage{amsmath,amssymb,amsthm,mathtools}
\usepackage{booktabs}
\usepackage{enumitem}
\usepackage[hidelinks]{hyperref}

\newtheorem{theorem}{Theorem}[section]
\newtheorem{lemma}[theorem]{Lemma}
\newtheorem{proposition}[theorem]{Proposition}
\newtheorem{corollary}[theorem]{Corollary}
\theoremstyle{definition}
\newtheorem{definition}[theorem]{Definition}
\newtheorem{remark}[theorem]{Remark}

\newcommand{\A}{A}
\newcommand{\C}{C}
\newcommand{\B}{B}
\newcommand{\Hh}{H}
\newcommand{\indeg}{d^-}
\newcommand{\outdeg}{d^+}
\newcommand{\lamax}{\lambda^{\max}}
\newcommand{\floor}[1]{\left\lfloor #1\right\rfloor}

\title{The simple directed connectivity problem at three paths}
\author{}
\date{}

\begin{document}
\begin{abstract}
For a simple digraph $D$, let $\lambda_D(u,v)$ be the maximum number of
arc-disjoint directed $u$--$v$ paths.  We determine the largest number of
arcs in an $n$-vertex simple digraph satisfying $\lambda_D(u,v)\le 2$ for
all ordered pairs.  The answer is
\[
 \max\left\{3(n-1),\floor{\frac{(n+1)^2}{4}}\right\}.
\]
The proof is human except for an exact finite verification through order
$10$.  Its main ingredients are a safe directed splitting operation and a
sharp counting lemma for posets in which every interval has at most three
elements.
\end{abstract}

\maketitle

\section{Statement and constructions}

For $n\ge 3$, write
\[
 g(n)=\max\bigl\{|A(D)|:\ |V(D)|=n,\ \lamax(D)\le 2\bigr\}
\]
and put
\[
 Q(n)=\floor{\frac{(n+1)^2}{4}},
 \qquad
 F(n)=\max\{3(n-1),Q(n)\}.
\]
A digraph is \emph{simple} here if it has no loops and at most one arc in
each direction on an unordered pair; digons are allowed.

\begin{theorem}\label{thm:main}
For every $n\ge 3$,
\[
 \boxed{g(n)=F(n)
 =\max\left\{3(n-1),\floor{\frac{(n+1)^2}{4}}\right\}.}
\]
\end{theorem}

Both lower bounds are immediate.  For the linear one, take a hub joined in
both directions to all other vertices and put a directed cycle on the
$n-1$ non-hub vertices.  There are $3(n-1)$ arcs, and every ordered pair
has one endpoint of in-degree or out-degree at most $2$.

For the quadratic one, split the vertices into sets $X,Y$, put every arc
from $X$ to $Y$, and put a directed cycle on $Y$.  Taking
$|Y|\in\{\lfloor(n+1)/2\rfloor,\lceil(n+1)/2\rceil\}$ gives
\[
 |Y|(|X|+1)=Q(n).
\]
A path ending at $y\in Y$ can use only the direct arc from its initial
vertex in $X$ or the unique cycle arc entering $y$; hence the maximum local
arc-connectivity is at most $2$.

It remains to prove the upper bound.  Notice that $Q(n)>3(n-1)$ from
$n=9$ onward.  Thus the induction below concerns the quadratic branch,
with the finite seam $3\le n\le 10$ treated exactly in
Section~\ref{sec:computer}.

\section{Safe splitting}

\begin{definition}
Let $v$ be a vertex of a simple digraph $D$.  A pair
\[
 x\to v,\qquad v\to y
\]
is \emph{legal} if $x\ne y$ and $x\to y\notin A(D)$.  Splitting this pair
means deleting the two displayed arcs and adding $x\to y$.

The \emph{split graph} at $v$ is the bipartite graph with left class
$N^-(v)$, right class $N^+(v)$, and an edge $xy$ precisely when the pair
$x\to v, v\to y$ is legal.  Let $\nu(v)$ be its matching number.
\end{definition}

\begin{lemma}[Safe splitting]\label{lem:splitting}
Let $M$ be a matching of size $r$ in the split graph at $v$.  Split all
pairs represented by $M$ and then delete $v$ and all its remaining incident
arcs.  The resulting simple digraph $D'$ satisfies
\[
 |A(D')|=|A(D)|-d(v)+r,
 \qquad
 \lamax(D')\le \lamax(D).
\]
\end{lemma}

\begin{proof}
The matching condition gives distinct incoming arcs and distinct outgoing
arcs at $v$.  The new arcs are pairwise distinct, are not loops, and were
absent before, so $D'$ is simple.

Take arc-disjoint directed paths in $D'$.  Replace every new arc $x\to y$
by the two-arc passage $x\to v\to y$ from which it arose.  Because the
split pairs form a matching, the lifted directed walks are still
arc-disjoint.  Erasing directed closed subwalks turns each of them into a
path without introducing an arc intersection.  Thus every family of
arc-disjoint paths in $D'$ lifts to one of the same size in $D$.
The arc count is immediate.
\end{proof}

The following consequence is the reason for introducing the operation.

\begin{corollary}\label{cor:split-bound}
Suppose Theorem~\ref{thm:main} is known for $n-1$, and let $D$ be an
$n$-vertex counterexample with exactly $F(n)+1$ arcs.  Then every vertex
$v$ satisfies
\[
 \nu(v)\le d(v)-\bigl(F(n)+1-F(n-1)\bigr).
\]
In particular, equality
$d(v)=F(n)+1-F(n-1)$ forces $\nu(v)=0$; equivalently,
\[
 x\to v,\ v\to y,\ x\ne y
 \quad\Longrightarrow\quad x\to y.
 \tag{\(*\)}\label{eq:transitive-vertex}
\]
\end{corollary}

\begin{proof}
Apply Lemma~\ref{lem:splitting} to a maximum matching and then use the
$(n-1)$-vertex upper bound.
\end{proof}

We call a vertex satisfying \eqref{eq:transitive-vertex}
\emph{transitive}.

\section{A poset lemma}

The low-degree vertices in a minimal even counterexample will all be
transitive.  Their induced digraph is therefore the comparability digraph
of a poset of height at most three, with a stronger interval restriction.
We isolate the exact estimate needed later.

\begin{lemma}[Rectangle packing]\label{lem:rectangles}
Let $A,C$ be sets of sizes $a,c$, and let
$X_i\times Y_i\subseteq A\times C$ $(1\le i\le b)$ be nonempty,
pairwise edge-disjoint rectangles.  Then
\[
 \sum_{i=1}^b (|X_i|+|Y_i|)
 \le \max\{a+bc,\ c+ab\}.
\]
If $c\ge a$, the right-hand side is $a+bc$.
\end{lemma}

\begin{proof}
Assume $c\ge a$ and put $x_i=|X_i|$, $y_i=|Y_i|$.
Edge-disjointness gives $\sum_i x_i y_i\le ac$.  Hence
\[
 c\Bigl(\sum_i x_i-a\Bigr)
 \le \sum_i x_i(c-y_i)
 \le c\sum_i(c-y_i).
\]
After cancelling $c$, this is
$\sum_i(x_i+y_i)\le a+bc$.  The other inequality follows by symmetry.
\end{proof}

\begin{lemma}[Posets with short intervals]\label{lem:poset}
Let $P$ be a finite poset on $N$ non-isolated elements such that every
interval contains at most three elements.  Let $A$ be its minimal elements,
$C$ its maximal elements, and $B$ the remaining elements, with respective
sizes $a,c,b$.  Then $B$ is an antichain.  For $z\in B$, put
\[
 X_z=\{x\in A:x<z\},
 \qquad
 Y_z=\{y\in C:z<y\}.
\]
The rectangles $X_z\times Y_z$ are nonempty and pairwise edge-disjoint.
After replacing $P$ by its dual if necessary, we may assume $c\ge a$, and
then
\begin{equation}\label{eq:poset-bound}
 e(P)\le c(N-c)+a.
\end{equation}
Moreover, the following near-equality statements hold.  Write
\[
 \delta=c(N-c)+a-e(P).
\]
\begin{enumerate}[label=\textup{(\roman*)},leftmargin=2.4em]
\item If $\delta=0$ and $c>a$, then all relations from $A$ to $C$ are
present, every $Y_z=C$, and the sets $X_z$ partition $A$.
\item If $\delta=0$, $c=a$, and $b\ge2$, the same conclusion holds after
possibly dualising $P$.
\item Suppose $b=2$, $\delta=1$, and $c-a\ge2$.  Then all relations from
$A$ to $C$ are present, the two sets $X_1,X_2$ are disjoint, and precisely
one of the following occurs:
\begin{enumerate}[label=\textup{(\alph*)}]
\item $Y_1=Y_2=C$ and $A\setminus(X_1\cup X_2)$ has one element;
\item $X_1\cup X_2=A$, one of $Y_1,Y_2$ equals $C$, and the other misses
one element of $C$.
\end{enumerate}
The same conclusion, after duality, holds when $c=a$.
\item If $b=1$, with middle element $z$, set
\[
 \alpha=a-|X_z|,
 \quad \beta=c-|Y_z|,
 \quad \gamma=ac-e(A,C).
\]
Then $\delta=\alpha+\beta+\gamma$; if $\delta>0$, then
$\alpha+\beta\ge1$.
\end{enumerate}
\end{lemma}

\begin{proof}
If $z<z'$ with $z,z'\in B$, choose $x<z$ minimal and $y>z'$ maximal.
Then $[x,y]$ has at least four elements, impossible.  Thus $B$ is an
antichain.  Every $X_z$ and $Y_z$ is nonempty.  If
$(x,y)\in(X_z\times Y_z)\cap(X_{z'}\times Y_{z'})$ for distinct $z,z'$,
then $[x,y]$ again contains four elements.  The rectangles are therefore
edge-disjoint.

Every comparability is either from $A$ to $C$, from $A$ to $B$, or from
$B$ to $C$, so
\[
 e(P)=e(A,C)+\sum_{z\in B}(|X_z|+|Y_z|).
\]
Equation~\eqref{eq:poset-bound} now follows from
$e(A,C)\le ac$ and Lemma~\ref{lem:rectangles}.

For the refinements, put
\[
 \gamma=ac-e(A,C),
 \qquad
 \rho=a+bc-\sum_{z\in B}(|X_z|+|Y_z|).
\]
Then $\delta=\gamma+\rho$, with $\gamma,\rho\ge0$.
Equality in the proof of Lemma~\ref{lem:rectangles} shows that if
$\rho=0$ and $c>a$, then every $Y_z=C$ and the $X_z$ partition $A$.
When $c=a$, rectangles of the two possible equality types cannot coexist,
so duality gives the same conclusion.  This proves (i) and (ii).

For (iii), if $X_1\cap X_2\ne\varnothing$, then
$Y_1\cap Y_2=\varnothing$, and consequently
\[
 \sum_{i=1}^2(|X_i|+|Y_i|)\le2a+c,
\]
which is at least $c-a\ge2$ below $a+2c$.  Hence $X_1,X_2$ are disjoint.
Now
\[
 \rho=\bigl(a-|X_1|-|X_2|\bigr)
       +(c-|Y_1|)+(c-|Y_2|).
\]
The case $\rho=0$ would force the rectangles to cover $A\times C$, and
then $\gamma=0$; thus $\delta=1$ forces $\rho=1$ and $\gamma=0$.
The two displayed alternatives are exactly the two ways in which three
nonnegative integers can sum to one.  If $c=a$, dualise first when the
$X_i$ are not disjoint.

Finally, when $b=1$, the identity in (iv) is immediate.  If
$\alpha=\beta=0$, then $X_z\times Y_z=A\times C$ forces
$e(A,C)=ac$, so also $\gamma=0$.
\end{proof}

For orientation arguments we use one more elementary obstruction.

\begin{lemma}[Saturated rectangle obstruction]\label{lem:saturated}
Let $L$ be a set of transitive vertices in a simple digraph with
$\lamax\le2$, and let $A_0,C_0\subseteq L$.  Suppose that for every
$a\in A_0$ there is $z(a)\in L$ such that, for every $c\in C_0$,
\[
 a\to c,
 \qquad a\to z(a),
 \qquad z(a)\to c.
 \tag{\(**\)}\label{eq:saturated}
\]
Let $H=V(D)\setminus L$.  Then for every $h\in H$,
\[
 h\not\to a,\qquad c\not\to h,
 \qquad\text{and}\qquad
 \{a\to h,\ h\to c\}\not\subseteq A(D)
\]
for all $a\in A_0$, $c\in C_0$.  In particular, if all arcs between
$H$ and $a,c$ are incident in the forced directions, then
\begin{equation}\label{eq:external-sum}
 d_H(a)+d_H(c)\le |H|.
\end{equation}
\end{lemma}

\begin{proof}
If $h\to a$, transitivity at $a$ and \eqref{eq:saturated} give
$h\to z(a)$ and $h\to c$.  The direct path $h\to c$ and the paths through
$a$ and $z(a)$ are three arc-disjoint $h$--$c$ paths.  Thus $h\not\to a$.
The assertion $c\not\to h$ follows dually, using transitivity at $c$.
Finally, $a\to h\to c$ cannot coexist with the direct path and the path
through $z(a)$.  Therefore the out-neighbours of $a$ in $H$ and the
in-neighbours of $c$ in $H$ are disjoint, proving
\eqref{eq:external-sum}.
\end{proof}

\section{Induction: the odd step}

We now prove the upper bound, assuming the finite verification of
Section~\ref{sec:computer}.

\begin{lemma}[Odd step]\label{lem:odd}
Let $k\ge5$.  If $g(2k)=k(k+1)$, then
\[
 g(2k+1)=(k+1)^2.
\]
\end{lemma}

\begin{proof}
The construction gives the lower bound.  Suppose an $(2k+1)$-vertex
feasible digraph has $(k+1)^2+1$ arcs.  Deleting any vertex leaves at most
$k(k+1)$ arcs, so every vertex has degree at least $k+2$.  But
\[
 \frac{2((k+1)^2+1)}{2k+1}<k+2
\]
for $k\ge5$, contradicting the average-degree bound.
\end{proof}

\section{Induction: the even step}

\begin{lemma}[Even step]\label{lem:even}
Let $k\ge6$.  If Theorem~\ref{thm:main} holds at every order below $2k$,
then
\[
 g(2k)=k(k+1).
\]
\end{lemma}

\begin{proof}
Suppose that a feasible digraph $D$ on $2k$ vertices has
\[
 E=k(k+1)+1
\]
arcs.  Since $g(2k-1)=k^2$, deletion gives $d(v)\ge k+1$ for every
vertex.  As
\[
 \sum_v d(v)=2E=2k(k+1)+2,
\]
exactly one of the following degree patterns occurs:
\begin{align*}
\text{(A)}&\quad (k+1)^{2k-1},\ k+3;\\
\text{(B)}&\quad (k+1)^{2k-2},\ (k+2)^2.
\end{align*}
Let $L$ be the vertices of degree $k+1$ and let $H=V(D)\setminus L$.
Corollary~\ref{cor:split-bound} gives $\nu(v)=0$ for every $v\in L$, so
all vertices of $L$ are transitive.

There is no digon inside $L$.  Indeed, deleting its two endpoints would
leave
\[
 E-2(k+1)+2=k(k-1)+1
\]
arcs on $2k-2$ vertices, whereas the induction hypothesis gives
$g(2k-2)=k(k-1)$.  Thus $D[L]$ is the comparability digraph of a poset
$P$.  Every interval of $P$ has at most three elements: two distinct
strict intermediate elements, together with the direct arc, would give
three arc-disjoint paths.  No vertex of $P$ is isolated, since a vertex of
$L$ can be incident with at most $2|H|\le4$ arcs outside $L$, while
$k+1\ge7$.

Let $A,B,C$ be the minimal, middle, and maximal levels from
Lemma~\ref{lem:poset}, with sizes $a,b,c$, after dualising so that $c\ge a$.
The set $B$ is nonempty.  Indeed, if $b=0$, then
$e(P)\le\floor{|L|^2/4}$, which is smaller than the value of $e(P)$ in
either degree pattern below.

\smallskip
\noindent\emph{Pattern (A): one exceptional vertex.}
Here $|H|=1$ and
\[
 N=|L|=2k-1,
 \qquad e(P)=E-(k+3)=k^2-2.
\]
Write
\[
 c=k-1+s,
 \qquad a=k-s-b.
\]
Lemma~\ref{lem:poset} gives
\[
 k^2-2\le c(N-c)+a=k^2-s^2-b,
\]
so $s^2+b\le2$.  The condition $c\ge a$ leaves exactly
\[
 (s,b)=(1,1),\ (0,1),\ (0,2).
\]

For $(1,1)$, the deficit in Lemma~\ref{lem:poset} is zero and the unique
middle vertex has internal degree $a+c=2k-2>k+1$, impossible.  For
$(0,1)$, the deficit is one; part (iv) of that lemma shows that the middle
vertex has internal degree at least $a+c-1=2k-3>k+1$.

It remains that $(s,b)=(0,2)$.  Then
$a=k-2$, $c=k-1$, and equality holds in
Lemma~\ref{lem:poset}.  Hence all $A\to C$ relations are present,
$Y_1=Y_2=C$, and $X_1,X_2$ partition $A$.  Each middle vertex has internal
degree $c+|X_i|\le k+1$, so $|X_i|\le2$.  Therefore $k=6$ and
$|X_1|=|X_2|=2$.  Every vertex of $A$ and $C$ then has internal degree
$6$ and exactly one incident arc to the exceptional vertex, whereas the
middle vertices have no external arc.  Conditions
\eqref{eq:saturated} hold with $A_0=A$, $C_0=C$ and the unique middle
vertex assigned to each $a$.  Equation~\eqref{eq:external-sum} gives
$1+1\le1$, a contradiction.

\smallskip
\noindent\emph{Pattern (B): two exceptional vertices.}
Let $r\in\{0,1,2\}$ be the number of arcs between the two vertices of $H$.
Then
\[
 N=|L|=2k-2,
 \qquad
 e(P)=E-\bigl(2(k+2)-r\bigr)=k(k-1)-3+r.
\]
Write
\[
 c=k-1+s,
 \qquad a=k-1-s-b.
\]
The poset bound becomes
\begin{equation}\label{eq:case-constraint}
 s^2+s+b\le3-r.
\end{equation}
Together with $c\ge a$, this gives precisely the cases in
Table~\ref{tab:cases}.  The symbol ``degree'' means that one middle vertex
already has internal degree greater than $k+1$.  The symbol
$(p,q)$ means that Lemma~\ref{lem:saturated} applies to nonempty sets
$A_0,C_0$ whose vertices have respectively $p$ and $q$ incident arcs to
$H$; then \eqref{eq:external-sum} says $p+q\le2$, a contradiction.

\begin{table}[ht]
\centering
\begin{tabular}{c c c c c}
\toprule
$r$ & $(s,b)$ & $(a,c)$ & $\delta$ & obstruction \\
\midrule
2 & $(0,1)$  & $(k-2,k-1)$ & 0 & degree \\
\midrule
1 & $(0,1)$  & $(k-2,k-1)$ & 1 & degree \\
1 & $(0,2)$  & $(k-3,k-1)$ & 0 & $(1,2)$ \\
1 & $(-1,2)$ & $(k-2,k-2)$ & 0 & $(2,1)$ \\
\midrule
0 & $(1,1)$  & $(k-3,k)$   & 0 & degree \\
0 & $(0,3)$  & $(k-4,k-1)$ & 0 & $(1,2)$ \\
0 & $(-1,3)$ & $(k-3,k-2)$ & 0 & $(2,1)$ \\
0 & $(0,2)$  & $(k-3,k-1)$ & 1 & $(1,2)$ \\
0 & $(-1,2)$ & $(k-2,k-2)$ & 1 & $(2,1)$ \\
0 & $(0,1)$  & $(k-2,k-1)$ & 2 & degree, except $k=6$ \\
\bottomrule
\end{tabular}
\caption{All integer solutions of \eqref{eq:case-constraint}.}
\label{tab:cases}
\end{table}

We justify the table entries.  In each zero-deficit row with $b\ge2$,
parts (i)--(ii) of Lemma~\ref{lem:poset} give, after duality when needed,
all $A\to C$ relations, $Y_z=C$, and a partition of $A$ by the sets $X_z$.
Consequently an $A$-vertex has internal degree $c+1$ and a $C$-vertex has
internal degree $a+b$.  Their external degrees are exactly the pair shown
in the final column, and Lemma~\ref{lem:saturated} applies.

The two one-deficit rows with $b=2$ follow from part (iii) of
Lemma~\ref{lem:poset}.  If one $A$-vertex is unused, take
$A_0=X_1\cup X_2$ and $C_0=C$.  If one middle vertex misses one maximal
element, take $A_0=A$ and delete that maximal element from $C_0$.  The
same external degrees, $(1,2)$ or $(2,1)$ after duality, hold on these
sets.

For the last row, let $z$ be the unique middle element and use the notation
$\alpha,\beta,\gamma$ from Lemma~\ref{lem:poset}(iv).  Its internal degree
is
\[
 a+c-(\alpha+\beta)\ge2k-5.
\]
This exceeds $k+1$ when $k>6$.  If $k=6$, equality forces
$\alpha+\beta=2$ and $\gamma=0$.  Put $A_0=X_z$, $C_0=Y_z$.
Every $a\in A_0$ has internal degree $6$ and external degree $1$, while
every $c\in C_0$ has internal degree $5$ and external degree $2$.
Lemma~\ref{lem:saturated} again gives $1+2\le2$, impossible.

All cases lead to a contradiction, proving the even step.
\end{proof}

\begin{proof}[Proof of Theorem~\ref{thm:main}]
The lower bounds were given in Section~1.  The finite verification in
Section~\ref{sec:computer} proves the theorem for $3\le n\le10$.
For $n\ge11$, use strong induction.  Lemma~\ref{lem:odd} proves every odd
order from the preceding even order, and Lemma~\ref{lem:even} proves every
even order at least $12$ from the smaller orders.
\end{proof}

\section{The finite verification}\label{sec:computer}

For completeness, we describe the exact finite calculation on which the
proof rests.  It checks that no feasible digraph has $F(n)+1$ arcs for
$3\le n\le10$; deleting arcs shows that this also excludes every larger
arc count.

There is one Boolean variable $a_{uv}$ for every possible arc $u\to v$.
For each ordered pair $(s,t)$, independent Boolean variables
$x_u^{st}$ choose a cut, with $x_s^{st}=1$ and $x_t^{st}=0$, and the model
imposes
\begin{equation}\label{eq:cut-encoding}
 \sum_{\substack{u\ne v\\x_u^{st}=1,\ x_v^{st}=0}} a_{uv}\le2.
\end{equation}
By directed max-flow/min-cut, the existence of such a cut for every
$(s,t)$ is exactly the condition $\lamax\le2$; this is not a relaxation.

For a target $K=F(n)+1$, deletion of one vertex gives
$d(v)\ge K-F(n-1)$.  The program enumerates all nondecreasing degree
sequences with this lower bound and sum $2K$.  It also inserts the valid
splitting inequalities from Corollary~\ref{cor:split-bound}: a vertex of
degree $d$ cannot have a legal-pair matching of size
$d-(K-F(n-1))+1$.  The remaining symmetry is broken by ordering
outdegrees inside equal-degree blocks.  At $n=9$, reversing every arc if
necessary lets one additionally assume that the first degree-$5$ vertex
has outdegree $0$, $1$, or $2$, and these three cases are checked
separately.

\begin{table}[ht]
\centering
\begin{tabular}{c c c c}
\toprule
$n$ & forbidden target $F(n)+1$ & degree sequences & verdict \\
\midrule
7 & 19 & 36 & all UNSAT \\
8 & 22 & 68 & all UNSAT \\
9 & 26 & 15 (45 reversal branches) & all UNSAT \\
10 & 31 & 2 & all UNSAT \\
\bottomrule
\end{tabular}
\caption{The nontrivial finite checks.  Orders $3$--$6$ are immediate
exhaustive cases.}
\label{tab:computer}
\end{table}

The accompanying file \texttt{verify\_m3\_bases.py} writes the exact
SMT-LIB instances and invokes Z3.  The source and the branch logs are part
of the supplementary material.  Thus the finite input to the induction is
fully reproducible.  The present verification records exact UNSAT solver
verdicts; it does not claim an independently replayable DRAT/LRAT proof
certificate.

\end{document}
